\documentclass[11pt,a4paper,hidelinks]{book}

%%%%%%%%%%%%%%%%%%%%%%%
%%% IMPORT PACKAGES %%%
%%%%%%%%%%%%%%%%%%%%%%%

% Use Unicode characters when possible
\usepackage[default,scale=0.95]{opensans}
\usepackage[utf8]{inputenc}
\usepackage[T1]{fontenc}

% adjust margins (should be done before importing fancyhdr!)
\usepackage[margin=2.5cm,a4paper]{geometry} %lmargin=2.5cm,rmargin=4cm,tmargin=3cm,bmargin=3cm

% text packages
\usepackage{csquotes}
\usepackage[english]{babel}
\usepackage{setspace}	

% scripting packages
\usepackage{ifthen}
\usepackage{blindtext}
\usepackage{todonotes}

% style packages
\usepackage{lmodern}
\usepackage{fancyhdr}
\usepackage[]{microtype} % Sublim­i­nal re­fine­ments to­wards ty­po­graph­i­cal per­fec­tion

% used for defining abbreviations, see nomenclature.tex
\usepackage{nomencl}

% figure packages
\usepackage{tikz}
\usetikzlibrary{calc,shapes,arrows,intersections,shadows}
\tikzset{>=latex}
\usepackage{graphicx}
\usepackage[indention=0.5cm,labelsep=colon,font={sf,small},labelfont={bf,sf}]{caption}
\usepackage[indention=0.5cm,font={sf,small},labelfont={bf,sf}]{subcaption}
\usepackage{float}
\usepackage{multirow}
\usepackage{import}
\usepackage{xcolor}
\usepackage{colortbl}

% pseudocode packages
\usepackage{algorithm}
\usepackage{algpseudocode}

% symbols packages
\usepackage{amssymb}
\usepackage{amsmath}

% table packages
\usepackage{tabularx}
\usepackage{tabulary}
\usepackage{booktabs}

% sections in allcaps
\usepackage{titlesec}
%\titleformat{\section} {\Large\bfseries\center}{}{0em}{\hrule height 0.4pt\vspace{2pt}}[\hrule height 0.4pt] \titlespacing*{\section} {0pt}{0pt}{10pt}
%\titleformat{\subsection} {\large\bfseries\center}{}{0em}{\vspace{0pt}} \titlespacing*{\subsection} {0pt}{0pt}{0pt}
%\titleformat{\subsubsection} {\normalsize\bfseries\center}{}{0em}{\vspace{0pt}} \titlespacing*{\subsubsection}  {0pt}{0pt}{0pt}
%\titleformat{\section} [block]{\Large\bfseries\filcenter}{}{0em}{\MakeUppercase}
%\titleformat{\subsection} [block]{\large\bfseries\filcenter}{}{0em}{\MakeUppercase}
%\titleformat{\subsubsection} [block]{\normalsize\bfseries\filcenter}{}{0em}{\MakeUppercase}

% references packages
\usepackage[sorting=none]{biblatex}
\bibliography{library}

% substitute in references
\DeclareDatafieldSet{setall}{
	\member[datatype=literal]
	\member[datatype=name]
	\member[field=journal]% journal is special since it is
	% actually journaltitle
}

\DeclareSourcemap{
	\maps[datatype=bibtex]{
		\map[overwrite, foreach={setall}]{
			% \`{\i}
			\step[fieldsource=\regexp{$MAPLOOP},
			match=\regexp{\x{0131}\x{0300}},
			replace=\regexp{\x{00EC}}]
			% \'{\i}
			\step[fieldsource=\regexp{$MAPLOOP},
			match=\regexp{\x{0131}\x{0301}},
			replace=\regexp{\x{00ED}}]
			% \^{\i}
			\step[fieldsource=\regexp{$MAPLOOP},
			match=\regexp{\x{0131}\x{0302}},
			replace=\regexp{\x{00EE}}]
			% \"{\i}
			\step[fieldsource=\regexp{$MAPLOOP},
			match=\regexp{\x{0131}\x{0308}},
			replace=\regexp{\x{00EF}}]
		}
	}
}

% link packages
\usepackage{url}
\usepackage{doi}

% add checkmark commands
\usepackage{pifont}
\newcommand{\cmark}{\ding{51}}
\newcommand{\xmark}{\ding{55}}

%%%%%%%%%%%%%%%%%%%%%
%%% DEFINE HEADER %%%
%%%%%%%%%%%%%%%%%%%%%
\usepackage{fancyhdr}  
\pagestyle{fancy} 
\renewcommand{\chaptermark}[1]{\markboth{#1}{}}
\fancypagestyle{frontmatter}{%    
	\fancyhf{} 	
	\fancyfoot[RO,LE]{\textsf{Page \thepage}} 
	\fancypagestyle{plain}{
		\renewcommand{\headrulewidth}{0pt}
		\fancyhead{}
		\fancyfoot[RO,LE]{\textsf{Page \thepage}}
	} 
}
\fancypagestyle{mainmatter}{%    
	\fancyhf{} 	
	\fancyhead[RO,LE]{\nouppercase{\small \textsf{\leftmark}}}
	\fancyfoot[RO,LE]{\textsf{Page \thepage}} 
}
\setlength{\headheight}{15pt}


\newcounter{letternum}
\newcounter{lettersum}
\setcounter{lettersum}{13}
\newlength{\thumbtopmargin}
\setlength{\thumbtopmargin}{1cm}
\newlength{\thumbbottommargin}
\setlength{\thumbbottommargin}{3cm}
\newlength{\thumbheight}
\pgfmathsetlength{\thumbheight}{%
	(\paperheight-\thumbtopmargin-\thumbbottommargin)/\value{lettersum}}

\newlength{\thumbwidth}
\setlength{\thumbwidth}{1.5cm}

\tikzset{
	thumb/.style={
		%   draw=black,
		fill=light-gray,
		text=black,
		minimum height=\thumbheight, %\thumbheight,
		text width=\thumbwidth,
		outer sep=0pt,%   outer sep=10pt,
		font=\sffamily\Huge,
	}
}
\newcommand{\oddthumb}[1]{%
	\begin{tikzpicture}[remember picture, overlay]
	\node [thumb,text centered,anchor=north east,] at ($%
	(current page.north east)-%
	(0,\thumbtopmargin+\value{letternum}*\thumbheight)%
	$) {#1};
	\end{tikzpicture}
}
\newcommand{\eventhumb}[1]{%
	\begin{tikzpicture}[remember picture, overlay]
	\node [thumb,text centered,anchor=north west,] at ($%
	(current page.north west)-%
	(0,\thumbtopmargin+\value{letternum}*\thumbheight)%
	$) {#1};
	\end{tikzpicture}
}
% create a new command to set a new lettergroup
\newcommand{\lettergroup}[1]{%
	\fancyhead[LO]{\oddthumb{#1}}%
	\fancyhead[RE]{\eventhumb{#1}}%
	\fancypagestyle{chapterstart}{%
		\renewcommand{\headrulewidth}{0pt}
		\renewcommand{\footrulewidth}{0pt}
		\fancyhf{}
		\chead{\oddthumb{#1}}% chapters start only on odd pages
		\fancyfoot[RO,LE]{\textsf{Page \thepage}} %\fancyfoot[RO,LE]{\textsf{Page \thepage}}
	}
	\thispagestyle{chapterstart}
	\stepcounter{letternum}%
}

%%%%%%%%%%%%%%%%%%%%%%%%%%%%%
%%% DEFINE FIGURE HELPERS %%%
%%%%%%%%%%%%%%%%%%%%%%%%%%%%%

\newcommand{\hugefigure}{\textwidth}
\newcommand{\LARGEfigure}{0.8\textwidth} % 1 per col
\newcommand{\LARgefigure}{0.6\textwidth} % 1 per col
\newcommand{\Largefigure}{0.5\textwidth} % 1 per col
\newcommand{\largefigure}{0.37\textwidth} % 2 per col
\newcommand{\normalfigure}{0.3\textwidth} % 3 per col
\newcommand{\smallfigure}{0.24\textwidth} % 3 per col
\newcommand{\Smallfigure}{0.19\textwidth} % 4 per col
\newcommand{\SMALLfigure}{0.15\textwidth} % 5 per col



\makeatletter
\newcommand*{\centerfloat}{%
	\parindent \z@
	\leftskip \z@ \@plus 1fil \@minus \textwidth
	\rightskip\leftskip
	\parfillskip \z@skip}
\makeatother

\newcommand{\picfig}[4]{
\begin{figure}[ht!]
	\centering \centerfloat
	\includegraphics[width=#2]{#1}
	\caption{#4}
	\label{#3}
\end{figure}}

%\newcommand{\picsubf}[2]{
%	\begin{subfigure}[b]{#2}
%		\fbox{\includegraphics[width=\linewidth]{#1}}
%\end{subfigure}}
%
%\newcommand{\picsubfc}[4]{
%	\begin{subfigure}[b]{#4}
%		\fbox{\includegraphics[width=\linewidth]{#1}}
%		\caption{#2}
%		\label{#3}
%\end{subfigure}}
%
%\newcommand{\picsub}[2]{
%	\begin{subfigure}[b]{#2}
%		\includegraphics[width=\linewidth]{#1}
%\end{subfigure}}
%
%\newcommand{\picsubc}[4]{
%	\begin{subfigure}[b]{#4}
%		\includegraphics[width=\linewidth]{#1}
%		\caption{#2}
%		\label{#3}
%\end{subfigure}}

%%%%%%%%%%%%%%%%%%%%%
%%% DEFINE STYLES %%%
%%%%%%%%%%%%%%%%%%%%%

% use Sonny chapter style
\usepackage[Sonny]{fncychap} %Sonny

\definecolor{light-gray}{gray}{0.90} %not always visible on a dell screen!!!

%% hyperlinks
% load hyperref after fncychap
\usepackage{hyperref}

% Set style variables
%\setlength{\parindent}{0cm} 
\renewcommand{\contentsname}{Table of Contents}
\renewcommand{\listfigurename}{List of Figures}
\renewcommand{\listtablename}{List of Tables}
\renewcommand{\appendixname}{}

\begin{document}

% thesis obligatory content: 
% * an abstract in Dutch
% * a list of abbreviations
% * one uniform (compiling) reference list
% * a resume or Curriculum Vitae of the candidate
% * an introductory chapter and additional linking texts between chapters (papers) if the dissertation consists of a compilation of papers
% * acknowledgements of financing institutions
% * apply for an ISBN number for your thesis. This takes 1-2 days.


\pagenumbering{alph}

\frontmatter
\pagestyle{empty}

\hypersetup{%
	pdftitle = {Modelling single cell dynamics with trajectories and gene regulatory networks},
	pdfsubject = {A subject},
	pdfkeywords = {import keywords},
	pdfauthor = {Robrecht Cannoodt},
	pdfcreator = {\LaTeX},
}

%%%%%%%%%%%%%%%%%%%%%%%%
%%% DEFINE TITLEPAGE %%%
%%%%%%%%%%%%%%%%%%%%%%%%
\begin{titlepage}
	
	\includegraphics[height=1.5cm]{fig/logos/logo_UGent_EN_RGB_color}
	
	\centering
	
	\vfill
	{\Huge \textsf{Modelling single cell dynamics\\with trajectories and gene regulatory networks}}
	
	\vspace{2.25cm}
	
	{\LARGE \textsf{Robrecht Cannoodt}}
	
	\vspace{2.25cm}
	
	{\textsf{Thesis submitted in fulfilment of the requirements for the degree of \\Doctor in Computer Science, 2019}}
	
	\vfill
	
	{\textsf{Supervisors: \\
			Prof. Dr. Yvan Saeys\\Prof. Dr. Katleen De Preter}}
	
	\vfill
	
	\includegraphics[height=.7cm]{fig/logos/icon_UGent_WE_EN_RGB_color}\hfill
	\includegraphics[height=1cm]{fig/logos/vib_rf_inflammation_research_rgb_pos} \hfill
	\includegraphics[height=.8cm]{fig/logos/dambi_logo} \hfill
	\includegraphics[height=.7cm]{fig/logos/cmgg_myshort} \hfill
	\includegraphics[height=.7cm]{fig/logos/FWO_Logo_KleurCMYK} 
	\vspace{1cm}
	
	{\textsf{Vakgroep Toegepaste Wiskunde, Informatica, en Statistiek\\
			Faculteit Wetenschappen, Universiteit Gent \\
			Krijgslaan 281 - S2, 9000 Gent}}
		
\end{titlepage}

%%%%%%%%%%%%%%%%%%%%%%%%
%%% ACKNOWLEDGEMENTS %%%
%%%%%%%%%%%%%%%%%%%%%%%%

\newpage{\thispagestyle{empty}\cleardoublepage}
\chapter*{Acknowledgements}

\newpage{\thispagestyle{empty}\cleardoublepage}

%%%%%%%%%%%%%
%%% LISTS %%%
%%%%%%%%%%%%%
\newpage{\thispagestyle{empty}\cleardoublepage}
\tableofcontents 

%\newpage{\thispagestyle{empty}\cleardoublepage}
%\listoffigures

%\newpage{\thispagestyle{empty}\cleardoublepage}
%\listoftables

\newpage{\thispagestyle{empty}\cleardoublepage}


%%%%%%%%%%%%%%%%%%%%
%%% NOMENCLATURE %%%
%%%%%%%%%%%%%%%%%%%%
\makenomenclature
\printnomenclature

\nomenclature{HCA}{Human Cell Atlas}
\nomenclature{GRN}{Gene Regulatory Network}
\nomenclature{IM}{Importance Measure}
\nomenclature{TF}{Transcription Factor}
\nomenclature{DNA}{Deoxyribonucleic Acid}
\nomenclature{RNA}{Ribonucleic Acid}
\nomenclature{mRNA}{Messenger RNA}
\nomenclature{NI}{Network Inference}
\nomenclature{ML}{Machine Learning}
\nomenclature{RF}{Random Forests}
\nomenclature{CART}{Classification And Regression Trees}


% chapters
\mainmatter
\pagestyle{mainmatter}

\newpage{\thispagestyle{empty}\cleardoublepage}
\chapter{Introduction} 
\chaptermark{Introduction.} % short title
\label{chap:introduction}
\lettergroup{\thechapter}

\begin{quote}
	\textbf{Abstract:} Recent developments in single-cell transcriptomics have opened new opportunities for studying dynamic processes in immunology in a high-throughput and unbiased manner. Starting from a mixture of cells in different stages of a developmental process, unsupervised trajectory inference algorithms aim to automatically reconstruct the underlying developmental path that cells are following. In this review, we break down the strategies used by this novel class of methods, and organize their components into a common framework, highlighting several practical advantages and disadvantages of the individual methods. We also give an overview of new insights these methods have already provided regarding the wiring and gene regulation of cell differentiation. As the trajectory inference field is still in its infancy, we propose several future developments which will ultimately lead to a global and data-driven way of studying immune cell differentiation.
\end{quote}

\vfill

Adapted from:\\
\textbf{Cannoodt, R.}$^*$, Saelens, W.$^*$, and Saeys, Y. Computational methods for trajectory inference from single-cell transcriptomics. \textit{European Journal of Immunology} 46, 11 (2016), 2496--2506. \doi{10.1002/eji.201646347}.\\
{\footnotesize $^*$ Equal contribution}
\newpage

\section{Single cell biology}

\subsection{Dynamic processes}

\subsection{Gene regulation}

\subsection{Single cell transcriptomics}

\section{Machine learning}

\subsection{Supervised learning}

\subsection{Unsupervised learning}

\subsection{Feature selection}

\subsection{Trajectory inference}

\subsection{Network inference}

\section{Research objectives}

\section{Outline}

% landmark mds
% princurve
% GillespieSSA
% dyngen

\newpage{\thispagestyle{empty}\cleardoublepage}
\chapter{A comparison of single-cell trajectory inference methods} 
\chaptermark{A comparison of single-cell trajectory inference methods.} % short title
\label{chap:dynbenchmark}

\begin{quote}
	\textbf{Abstract:} \blindtext
\end{quote}

\vfill

Adapted from:\\
Saelens, W.$^*$, \textbf{Cannoodt, R.}$^*$, Todorov, H., and Saeys, Y. A comparison of single-cell trajectory inference methods. \textit{Nature Biotechnology} 37, 5 (2019), 547--554. \doi{10.1038/s41587-019-0071-9}.\\
{\footnotesize $^*$ Equal contribution}
\newpage

\blindtext



\newpage{\thispagestyle{empty}\cleardoublepage}
\chapter{SCORPIUS} 
\chaptermark{SCORPIUS.} % short title
\label{chap:scorpius}

\begin{quote}
	\textbf{Abstract:} \blindtext
\end{quote}

\vfill

Adapted from:\\
\textbf{Cannoodt, R.}, \ldots, De Preter, K., and Saeys, Y. SCORPIUS improves trajectory inference and identifies novel modules in dendritic cell development. \textit{Journal} vol, issue (2019), page--page. \doi{10.1101/079509v2}.
\newpage

\blindtext

% include gng ?

\newpage{\thispagestyle{empty}\cleardoublepage}
\chapter{dyno: A toolkit for inferring and interpreting trajectories} 
\chaptermark{dyno: A toolkit for inferring and interpreting trajectories.} % short title
\label{chap:dyno}

\begin{quote}
	\textbf{Abstract:} \blindtext
\end{quote}

\vfill

Adapted from:\\
\textbf{Cannoodt, R.}, Saelens, W., and Saeys, Y. dyno: A toolkit for inferring and interpreting trajectories. \textit{Journal} vol, issue (2019), page--page. doi.
\newpage

\blindtext

\newpage{\thispagestyle{empty}\cleardoublepage}
\chapter{Inferring single cell regulatory networks} 
\chaptermark{Inferring single cell regulatory networks.} % short title
\label{chap:brood}

\begin{quote}
	\textbf{Abstract:} \blindtext
\end{quote}

\vfill

Adapted from:\\
\textbf{Cannoodt, R.}, Saelens, W., and Saeys, Y. Inferring Single Cell Regulatory Networks. \textit{Journal} vol, issue (2019), page--page. doi.
\newpage

\blindtext

\newpage{\thispagestyle{empty}\cleardoublepage}
\chapter{Optimising regulatory networks} 
\chaptermark{Optimising regulatory networks.} % short title
\label{chap:incgraph}

\begin{quote}
	\textbf{Abstract:} \blindtext
\end{quote}

\vfill

Adapted from:\\
\textbf{Cannoodt, R.}, Ruyssinck, J., Ramon, J., De Preter, K., and Saeys, Y. IncGraph: Incremental graphlet counting for topology optimisation. \textit{PLOS ONE} 13, 4 (2018), e0195997. \doi{10.1371/journal.pone.0195997}.
\newpage

\blindtext \cite{Eren2013}

\newpage{\thispagestyle{empty}\cleardoublepage}
\chapter{General discussion} 
\chaptermark{General discussion.} % short title
\label{chap:discussion}

\blindtext

\appendix


\backmatter


\chapter{Samenvatting}

\chapter{Summary}

\chapter{List of Publications}

\printbibliography

% thanking:
% the people who brought me soup and sandwiches (look up names)
% the people who tried to use the PRISM cluster but couldn't


\end{document}
